% GENERATED FILE. Edit canonical lessons, then run npm run generate:books.
% canonical-source-hash: fnv1a64:b60313411601d440
% canonical-lessons: GU-C491-jamai, GU-C491-salo, GU-C491-banevi, GU-C491-vevai, GU-C491-vevan

\chapter{A Wedding in the Family}
\label{ch:gu-a-wedding-in-the-family}
\hlchaptermodality{491}

\begin{chapteropening}
\textbf{By the end of this chapter:} \emph{I can name relatives by marriage.}

Complete the last lesson of chapter 491: I can name relatives by marriage.
\end{chapteropening}

\section[\texorpdfstring{\gu{જમાઈ} (jamāī)}{જમાઈ (jamāī)}]{\gu{જમાઈ} (jamāī) — a son-in-law}

\label{lesson:GU-C491-jamai}



\begin{quote}
Before the new one: say the Gujarati for a husband's elder brother's wife, then the Gujarati for married life.
\end{quote}

\subsection*{You'll want to know: \gu{જમાઈ}}
\textbf{\gu{જમાઈ}} — \emph{jamāī} — \textquotedblleft{}a son-in-law\textquotedblright{}.

\begin{quote}
\gu{અમારા} \gu{જમાઈ} \gu{ઘણા} \gu{નમ્ર} \gu{છે}.
\end{quote}

\begin{quote}
\emph{amārā jamāī ghaṇā namra che.}
\end{quote}

\begin{quote}
Our son-in-law is very polite.
\end{quote}

A masculine noun.

\begin{cousinweb}
\textbf{\gu{જમાઈ}} (\emph{jamāī}) is inherited from Sanskrit \emph{jāmātṛ}, \textquotedblleft{}daughter's husband\textquotedblright{}; CDIAL lists the Gujarati word (CDIAL 5198).
\end{cousinweb}

\begin{grammarlens}[title={What you've built}]
One more word for describing what you lived and felt.
\end{grammarlens}

\subsection*{Your turn}
\begin{itemize}
\raggedright
  \item \emph{Say it:} \emph{jamāī}
  \item \emph{Say it:} \emph{jamāī}, once more
  \item \emph{Say it:} \emph{amārā jamāī ghaṇā namra che.} — Our son-in-law is very polite.
  \item \emph{Recall:} say the Gujarati for a husband's elder brother's wife, then the Gujarati for married life, then say \emph{jamāī} again
\end{itemize}

\begin{tcolorbox}[breakable,colback=teal!4,colframe=teal!35!black,title={Before you move on}]
What does \gu{જમાઈ} mean? (\textquotedblleft{}A son-in-law\textquotedblright{}.) Say it once more.
\end{tcolorbox}
\section[\texorpdfstring{\gu{સાળો} (sāḷo)}{સાળો (sāḷo)}]{\gu{સાળો} (sāḷo) — a wife's brother}

\label{lesson:GU-C491-salo}



\begin{quote}
Before the new one: say the Gujarati for married life, then the Gujarati for a son-in-law.
\end{quote}

\subsection*{You'll want to know: \gu{સાળો}}
\textbf{\gu{સાળો}} — \emph{sāḷo} — \textquotedblleft{}a wife's brother\textquotedblright{}.

\begin{quote}
\gu{અરુણનો} \gu{સાળો} \gu{શહેરમાં} \gu{નોકરી} \gu{કરે} \gu{છે}.
\end{quote}

\begin{quote}
\emph{Aruṇno sāḷo śahermā̃ nokrī kare che.}
\end{quote}

\begin{quote}
Arun's wife's brother works in the city.
\end{quote}

A masculine noun.

\begin{cousinweb}
\textbf{\gu{સાળો}} (\emph{sāḷo}) is inherited from Sanskrit \emph{syāla}, \textquotedblleft{}wife's brother\textquotedblright{}; CDIAL lists the Gujarati word (CDIAL 13871).
\end{cousinweb}

\begin{grammarlens}[title={What you've built}]
One more word for describing what you lived and felt.
\end{grammarlens}

\subsection*{Your turn}
\begin{itemize}
\raggedright
  \item \emph{Say it:} \emph{sāḷo}
  \item \emph{Say it:} \emph{sāḷo}, once more
  \item \emph{Say it:} \emph{Aruṇno sāḷo śahermā̃ nokrī kare che.} — Arun's wife's brother works in the city.
  \item \emph{Recall:} say the Gujarati for married life, then the Gujarati for a son-in-law, then say \emph{sāḷo} again
\end{itemize}

\begin{tcolorbox}[breakable,colback=teal!4,colframe=teal!35!black,title={Before you move on}]
What does \gu{સાળો} mean? (\textquotedblleft{}A wife's brother\textquotedblright{}.) Say it once more.
\end{tcolorbox}
\section[\texorpdfstring{\gu{બનેવી} (banevī)}{બનેવી (banevī)}]{\gu{બનેવી} (banevī) — a sister's husband}

\label{lesson:GU-C491-banevi}



\begin{quote}
Before the new one: say the Gujarati for a son-in-law, then the Gujarati for a wife's brother.
\end{quote}

\subsection*{You'll want to know: \gu{બનેવી}}
\textbf{\gu{બનેવી}} — \emph{banevī} — \textquotedblleft{}a sister's husband\textquotedblright{}.

\begin{quote}
\gu{મારા} \gu{બનેવીએ} \gu{મને} \gu{જન્મદિવસે} \gu{ઘડિયાળ} \gu{આપી}.
\end{quote}

\begin{quote}
\emph{mārā banevīe mane janmadivse ghaḍiyāḷ āpī.}
\end{quote}

\begin{quote}
My sister's husband gave me a watch for my birthday.
\end{quote}

A masculine noun.

\begin{cousinweb}
\textbf{\gu{બનેવી}} (\emph{banevī}) is inherited from Sanskrit \emph{bhaginīpati}, \textquotedblleft{}sister's husband\textquotedblright{}, from \emph{bhaginī}, \textquotedblleft{}sister\textquotedblright{}, and \emph{pati}, \textquotedblleft{}husband\textquotedblright{}; CDIAL lists the Gujarati word (CDIAL 9350).
\end{cousinweb}

\begin{grammarlens}[title={What you've built}]
One more word for describing what you lived and felt.
\end{grammarlens}

\subsection*{Your turn}
\begin{itemize}
\raggedright
  \item \emph{Say it:} \emph{banevī}
  \item \emph{Say it:} \emph{banevī}, once more
  \item \emph{Say it:} \emph{mārā banevīe mane janmadivse ghaḍiyāḷ āpī.} — My sister's husband gave me a watch for my birthday.
  \item \emph{Recall:} say the Gujarati for a son-in-law, then the Gujarati for a wife's brother, then say \emph{banevī} again
\end{itemize}

\begin{tcolorbox}[breakable,colback=teal!4,colframe=teal!35!black,title={Before you move on}]
What does \gu{બનેવી} mean? (\textquotedblleft{}A sister's husband\textquotedblright{}.) Say it once more.
\end{tcolorbox}
\section[\texorpdfstring{\gu{વેવાઈ} (vevāī)}{વેવાઈ (vevāī)}]{\gu{વેવાઈ} (vevāī) — the father-in-law of one's child}

\label{lesson:GU-C491-vevai}



\begin{quote}
Before the new one: say the Gujarati for a wife's brother, then the Gujarati for a sister's husband.
\end{quote}

\subsection*{You'll want to know: \gu{વેવાઈ}}
\textbf{\gu{વેવાઈ}} — \emph{vevāī} — \textquotedblleft{}the father-in-law of one's child\textquotedblright{}.

\begin{quote}
\gu{લગ્નમાં} \gu{બંને} \gu{વેવાઈ} \gu{ખુશ} \gu{હતા}.
\end{quote}

\begin{quote}
\emph{lagnamā̃ banne vevāī khuś hatā.}
\end{quote}

\begin{quote}
At the wedding both fathers-in-law were happy.
\end{quote}

A masculine noun.

\begin{cousinweb}
\textbf{\gu{વેવાઈ}} (\emph{vevāī}) is inherited from Sanskrit \emph{vaivāhika}, \textquotedblleft{}nuptial, related by marriage\textquotedblright{}; CDIAL lists the Gujarati word, \textquotedblleft{}son's or daughter's father-in-law\textquotedblright{} (CDIAL 12148).
\end{cousinweb}

\begin{grammarlens}[title={What you've built}]
One more word for describing what you lived and felt.
\end{grammarlens}

\subsection*{Your turn}
\begin{itemize}
\raggedright
  \item \emph{Say it:} \emph{vevāī}
  \item \emph{Say it:} \emph{vevāī}, once more
  \item \emph{Say it:} \emph{lagnamā̃ banne vevāī khuś hatā.} — At the wedding both fathers-in-law were happy.
  \item \emph{Recall:} say the Gujarati for a wife's brother, then the Gujarati for a sister's husband, then say \emph{vevāī} again
\end{itemize}

\begin{tcolorbox}[breakable,colback=teal!4,colframe=teal!35!black,title={Before you move on}]
What does \gu{વેવાઈ} mean? (\textquotedblleft{}The father-in-law of one's child\textquotedblright{}.) Say it once more.
\end{tcolorbox}
\section[\texorpdfstring{\gu{વેવાણ} (vevāṇ)}{વેવાણ (vevāṇ)}]{\gu{વેવાણ} (vevāṇ) — the mother-in-law of one's child}

\label{lesson:GU-C491-vevan}



\begin{quote}
Before the new one: say the Gujarati for a sister's husband, then the Gujarati for the father-in-law of one's child.
\end{quote}

\subsection*{You'll want to know: \gu{વેવાણ}}
\textbf{\gu{વેવાણ}} — \emph{vevāṇ} — \textquotedblleft{}the mother-in-law of one's child\textquotedblright{}.

\begin{quote}
\gu{વેવાણે} \gu{મહેમાનોનું} \gu{ઉમળકાથી} \gu{સ્વાગત} \gu{કર્યું}.
\end{quote}

\begin{quote}
\emph{vevāṇe mahemānonũ umaḷkāthī svāgat karyũ.}
\end{quote}

\begin{quote}
Our child's mother-in-law welcomed the guests warmly.
\end{quote}

A feminine noun.

\begin{cousinweb}
\textbf{\gu{વેવાણ}} (\emph{vevāṇ}) is inherited from Sanskrit \emph{vivāhin}, \textquotedblleft{}connected by marriage\textquotedblright{}; CDIAL lists the Gujarati word, \textquotedblleft{}son's or daughter's mother-in-law\textquotedblright{} (CDIAL 11925).
\end{cousinweb}

\begin{grammarlens}[title={What you've built}]
One more word for describing what you lived and felt.
\end{grammarlens}

\subsection*{Your turn}
\begin{itemize}
\raggedright
  \item \emph{Say it:} \emph{vevāṇ}
  \item \emph{Say it:} \emph{vevāṇ}, once more
  \item \emph{Say it:} \emph{vevāṇe mahemānonũ umaḷkāthī svāgat karyũ.} — Our child's mother-in-law welcomed the guests warmly.
  \item \emph{Recall:} say the Gujarati for a sister's husband, then the Gujarati for the father-in-law of one's child, then say \emph{vevāṇ} again
\end{itemize}

\begin{tcolorbox}[breakable,colback=teal!4,colframe=teal!35!black,title={Before you move on}]
What does \gu{વેવાણ} mean? (\textquotedblleft{}The mother-in-law of one's child\textquotedblright{}.) Say it once more.
\end{tcolorbox}
